OpenFisca-France (législation des revenus 2024) ; DGFiP, IRCOM revenus 2024 ; calculs de l'auteur (synthese/modele\_purs.py, architectures\_ircsg.py, etapes\_ircsg.py).\par Note : prélèvement sur 1 000 \euro{} de revenu brut en plus (IR + CSG + CRDS, et Casa pour le retraité), cas-types sans niche ; lissé sur 2 000 \euro{}. Étape 1 : barème de l'IR (0, 11, 30, 41, 45 \%) appliqué à chaque adulte sur le revenu après cotisations, CSG et CRDS remplacées par un taux unique de 1,55 \%. Étape 2 : taux marginal en ligne droite, le barème le plus proche de ce que chacun paie aujourd'hui. Proposition : franchise de 4 000 \euro{} sur les revenus d'activité, pensions exonérées jusqu'à 12 500 \euro{} puis de moins en moins jusqu'à 25 000 \euro{} (d'où le taux de 28 \% du retraité dans cette zone).
